\section{Binäre Zahlendarstellung von Festkommazahlen}

Im Rahmen dieses Projekts wird von Ein- sowie Ausgangswerten mit einer Genauigkeit von 12 Bit ausgegangen.
Basierend auf älteren Sensoren wird von Werten im Bereich von $-2 < z < 2$ ausgegangen.
Aus diesem Grund müssen sowohl ein Ganzzahlanteil, sowie Nachkommastellen repräsentiert werden können. Wie dies gelingt, wird in den nächsten Abschnitten gezeigt.
Hierfür werden Festkommazahlen verwendet, aufgrund der Rechenoperationen haben diese dennoch unterschiedlich viele Vor- sowie Nachkommastellen.


\subsection{Integer-Zahl im 1er-Komplement}
Bei der Interpretation des Bitvektors als Integerwert im Einerkomplement werden die Bits anhand ihrer Position im Bitvektor gewichtet, wobei das niederwertigste Bit 
(LSB, least significant bit) dem Wert für den Faktor $2^0$ entspricht, das Bit links davon dem für $2^1$ und so weiter. Die Summe aller Bits, ohne das höchstwertigste, 
multipliziert mit ihrer Wertigkeit (Potenz) ergibt den Betrag der Dezimalzahl. Das höchstwertigste Bit (MSB, most significant bit) gibt Auskunft darüber, ob es sich 
um eine negative oder positive Zahl handelt, wobei eine 0 für eine positive Zahl steht. Entsprechend besagt die 1, dass die Zahl negativ ist.
Dies hat zur Folge, dass es eine positive und eine negative Null und somit eine Doppeldeutigkeit gibt. Des Weiteren wird
ein LSB an Auflösung verschenkt. Der Wertebereich erstreckt sich von $-2^{\mathrm{MSB}-1}+1\,\mathrm{LSB}$ bis $2^{\mathrm{MSB}-1}-1\,\mathrm{LSB}$

Diese Darstellung hat den Vorteil, dass sich das Ergebnis einer Multiplikation der Zahlen $a \cdot b$ und $-a \cdot b$ nur im vordersten Bit unterscheidet. Darüber hinaus
lässt sich das Vorzeichen des Ergebnisses durch eine einfache XOR-Verknüpfung der beiden MSB der Multiplikanden ermitteln. 
Die eigentliche Multiplikation beschränkt sich auf die Bits MSB-1 bis LSB.

Nachteile zeigen sich hingegen bei der Addition sowie Subtraktion negativer Zahlen. Auch hierfür gibt es schematische Rechenregeln, diese erfordern jedoch mehr 
Zwischenschritte als im Zweierkomplement. 


\subsection{Integer-Zahl im 2er-Komplement}\label{sec:Integer2erKomplement}

Bei der Interpretation als Zweierkomplement kann anhand es MSB ebenfalls erkannt werden, ob es sich um eine positive oder negative Zahl handelt. 
Hie bedeutet ein gesetztes MSB $-2^{MSB-1}$, was der negativsten darstellbaren Zahl entspricht. Hierbei sind alle anderen 
Bits auf 0. Für gesetzte Bits wird der Dezimalwert, wie beim Einerkomplement beschrieben, berechnet und auf den negativen Wert aufaddiert. Wenn das MSB nicht gesetzt
ist, wird der errechnete Dezimalwert auf 0 addiert. Auf diese Weise lassen sich Zahlen im Wertebereich von $-2^{MSB-1}$ bis $2^{MSB-1}-1 \,LSB$ darstellen. Der positive
Wertebereich ist also um ein LSB kleiner als der negative und es gibt keine doppelte Null.

Um das Vorzeichen umzukehren müssen alle Bits invertiert werden. Auf das Resultat muss abschließend 1 LSB addiert werden.

Vorteil bei dieser Darstellung ist, dass die mathematischen Operationen Addition, Subtraktion und Multiplikation direkt angewandt werden können. Unterstützt werden sie z.B. 
von den Datentypen \texttt{unsigned} sowie \texttt{signed}, welche in der Bibliothek u.a. \texttt{ieee.numeric\_std.all} definiert sind.


\subsection{Darstellung dualer Zahlen im SQ-Format}
Im SQ-Format werden Zahlen als vorzeichenbehafteter Quotient (signed quotient) dargestellt. Wie beim 2er-Komplement entscheidet das höchstwertigste Bit, ob es sich um eine
positive oder negative Zahl handelt. In Abbildung \ref{pic:SQKreis} ist exemplarisch die Interpretation von Dualzahlen im SQ3-Format, also für vier Bit, zu sehen.

\begin{figure}[ht!]
 \centering
 \includegraphics[width=0.5\textwidth]{img/SQ-Kreis.png}
 \caption{Interpretation von Dualzahlen im SQ3-Format.}
 \label{pic:SQKreis}
\end{figure}


Der darstellbare Zahlenbereich liegt hier bei $-1\leq z < 1$. Benötigt werden Zahlen im Bereich von etwa $\pm2$, weshalb ein Vorkommabit benötigt wird. 
Da 12 Bit zur Verfügung stehen, von denen eins für das Vorzeichen und ein weiteres für eine Vorkommastelle verwendet werden, bleiben 10 Bits für die Nachkommazahlen übrig.
Die Aufteilung der Bits wird über die Bezeichnung S1Q10 definiert.
Da für den Quotient 10 Bit zur Verfügung stehen, beträgt die maximale Auflösung $1\,LSB = 2^{-10} = {1024}^{-1} = 9,765625\cdot10^{-4}$.
Der Wertebereich liegt in diesem Fall liegt bei \num{-2} bis \num{1,999023438}. 

Für die Addition oder Multiplikation zweier Zahlen müssen beide einerseits dieselbe Bitbreite und andererseits das gleiche Darstellungsformat besitzen.



 
\subsection{Numerisch bedingte Ungenauigkeiten}\label{sec:NumerischeUngenauigkeiten}
Numerische Ungenauigkeiten entstehen immer dann, wenn die zur Verfügung stehenden Bits es nicht ermöglichen eine Zahl exakt abzubilden. 
Bei einem Bitshift, welcher häufig für die Division durch Zwei oder Vielfachen von Zwei verwendet wird, kann immer Information verloren gehen. Dies ist immer dann der Fall,
wenn die Bits die abgeschnitten werden eine 1 sind. Das hat zur Folge, dass beispielsweise
bei einer Division durch Zwei der resultierende Wert um 1\,LSB kleiner ist, als er eigentlich sein sollte. 
Dieses Problem kann bei jedem Bitshift auftreten. Die Wahrscheinlichkeit für eine 1 liegt im Mittel bei \SI{50}{\percent}, weshalb davon ausgegangen werden muss, dass ein 
positives Ergebnis etwas kleiner und ein negatives vom Betrag her etwas größer ist, als bei verlustfreier Berechnung. 

% Von Prof. Vollmer: 
% \begin{equation}
%  S_N = \frac{P_Q}{\left(2^0\right)^2} + \dfrac{P_Q}{\left(2^1\right)^2} + \dfrac{P_Q}{\left(2^2\right)^2} + \cdots + \dfrac{P_Q}{\left(2^{L-1}\right)^2}
% \end{equation}
% L : Stufe (Additionstakte?)

Da diese Arbeit den Schwerpunkt in der Aufwandsabschätzung einer Chipimplementation einer 2D-DFT auf einem \gls{asic} hat, ist diese Problematik kein Gegenstand dieser Arbeit und
wird an dieser Stelle nur in Grundzügen erwähnt. %Für eine 

